%Chargement des paquets
\usepackage{amsmath}
\usepackage{amsfonts}
\usepackage{amssymb}
\usepackage{enumerate}
\usepackage{mathtools}
\usepackage{bbm}
\usepackage{xparse, etoolbox}
\usepackage{enumerate}
\usepackage{enumitem}
\usepackage{mathabx}
\usepackage{babel}[french]

%environment
\usepackage{amsthm}
\usepackage{keytheorems}
\usepackage{hyperref} 

%Interface théorème
\newkeytheorem{lem}[
	name = Lemme
]

\newkeytheorem{prop}[
	name = Proposition
]

\newkeytheorem{thm}[
	name = Théorème
]

\newkeytheorem{coro}[
	name = Corollaire
]

\renewcommand*{\proofname}{Démonstration}

\theoremstyle{definition}
\newtheorem{defn}{Définition}

\theoremstyle{definition}
\newtheorem*{nota}{Notation}

\theoremstyle{definition}
\newtheorem*{limi}{Liminaire}

\theoremstyle{remark}
\newtheorem{rmq}{Remarque}

\theoremstyle{plain}
\newtheorem*{fait}{Fait}


%Ensembles de nombres
\newcommand{\N} {\mathbb{N}}
\newcommand{\Ne}{\N^\ast}
\newcommand{\Z} {\mathbb{Z}}
\newcommand{\Q} {\mathbb{Q}}
\newcommand{\R} {\mathbb{R}}
\newcommand{\Rb}{\overline{\mathbb{R}}}
\newcommand{\Rp}{\R_+}
\newcommand{\Rm}{\R_-}
\newcommand{\K} {\mathbb{K}}
\newcommand{\Ket} {\mathbb{K^{\ast}}}
\newcommand{\Cx}{\mathbb{C}}

%Ensembles, tribus
\newcommand{\A}{\mathbb{A}}
\renewcommand{\P}{\mathcal{P}}
\newcommand{\T}{\mathcal{T}}
\newcommand{\F}{\mathcal{F}}
\newcommand{\C} {\mathcal{C}}
\newcommand{\ouv}{\mathcal{O}}
\newcommand{\B}{\mathcal{B}}
\newcommand{\M}{\mathcal{M}}
\newcommand{\etag}{\mathcal{E}}
\newcommand{\etagOp}{\etag^{+}(\Omega)}
\newcommand{\etagO}{\etag(\Omega)}
%%Lp avant quotientage
\newcommand{\Lic}{\mathcal{L}^1}
\newcommand{\Lpc}{\mathcal{L}^p}
\newcommand{\Linfc}{\mathcal{L}^{\infty}}
\newcommand{\Lifull}{\Lic(\Omega, \B, m)}
\newcommand{\Lpfull}{\Lpc(\Omega, \B, m)}
\newcommand{\Linffull}{\Linfc(\Omega, \B, m)}
%%Lp après quotientage
\newcommand{\Li}{L^1}
\newcommand{\Ld}{L^2}
\newcommand{\Lp}{L^p}
\newcommand{\Linf}{L^{\infty}}
\newcommand{\Listd}{\Li(\Omega, m)} 
\newcommand{\Ldstd}{\Ld(\Omega, m)} 
\newcommand{\Lpstd}{\Lp(\Omega, m)} 

%Quelques opérateurs
\DeclareMathOperator{\codim}{codim}
\DeclareMathOperator{\rg}{rg}
\DeclareMathOperator{\tr}{tr}
\DeclareMathOperator{\vect}{Vect}
\DeclareMathOperator{\ima}{Im}
\DeclareMathOperator{\id}{Id}
\DeclareMathOperator{\sign}{sign}
\DeclareMathOperator{\ext}{ext}
\newcommand{\ind}{\mathbbm{1}}
\newcommand*{\spr}[2]{\left(\,#1\,\middle|\,#2\,\right)}
\newcommand{\interior}[1]{%
  {\kern0pt#1}^{\mathrm{o}}%
}
\DeclareMathOperator{\card}{card}
\DeclareMathOperator{\ppcm}{ppcm}

%Commandes de pure flemme
\newcommand{\veps}{\varepsilon}
\newcommand{\intab}{\int_{a}^{b}}
\newcommand{\intR}{\int_{-\infty}^{\infty}}
\newcommand{\intinf}{\int_{- \infty}^{+ \infty}}
\newcommand{\equi}{\Leftrightarrow}
\newcommand{\Kev}{$\K$-espace vectoriel }
\newcommand{\Kevs}{$\K$-espaces vectoriels }
\newcommand{\Rev}{$\R$-espace vectoriel }
\newcommand{\Revs}{$\R$-espaces vectoriels }
\newcommand{\Cev}{$\Cx$-espace vectoriel }
\newcommand{\Cevs}{$\Cx$-espaces vectoriels }
\newcommand{\evn}{(E, \norm{.})}
\newcommand{\interf}[2]{[#1,#2]}
\newcommand{\limb}{\lim\limits_}
\newcommand{\lebn}{\lambda_n}
\newcommand{\pp}{\text{~presque partout}}
\newcommand{\derivp}[2]{\frac{\partial #1}{\partial #2}}
\newcommand{\famiu}{(u_i)_{i \in I}}
\newcommand{\sev}{\text{~s.e.v.}}
\newcommand{\Mnp}{M_{n,p}(\K)}
\newcommand{\Mn}{M_{n}(\K)}
\newcommand{\tq}{\text{~t.q.~}}
\newcommand{\mq}{~montrer~que~}
\newcommand{\et}{\quad \text{et} \quad}
\newcommand{\ou}{\text{~ou~}}
\newcommand{\Kx}{\K[X]}


%Commandes de gars chelou
\renewcommand{\subset}{\subseteq}

%Commande restriction, ty duckduckgo
\def\restr#1#2{\mathchoice
              {\setbox1\hbox{${\displaystyle #1}_{\scriptstyle #2}$}
              \restraux{#1}{#2}}
              {\setbox1\hbox{${\textstyle #1}_{\scriptstyle #2}$}
              \restraux{#1}{#2}}
              {\setbox1\hbox{${\scriptstyle #1}_{\scriptscriptstyle #2}$}
              \restraux{#1}{#2}}
              {\setbox1\hbox{${\scriptscriptstyle #1}_{\scriptscriptstyle #2}$}
              \restraux{#1}{#2}}}
\def\restraux#1#2{{#1\,\smash{\vrule height .8\ht1 depth .85\dp1}}_{\,#2}} 

%Quelques normes
\DeclarePairedDelimiter\abs{\lvert}{\rvert}
\DeclarePairedDelimiter\labs{\bigl\lvert}{\bigr\rvert}
\DeclarePairedDelimiter\norm{\lVert}{\rVert}
\DeclarePairedDelimiterXPP\normi[1]{}\lVert\rVert{_1}{\ifblank{#1}{\:\cdot\:}{#1}}
\DeclarePairedDelimiterXPP\normd[1]{}\lVert\rVert{_2}{\ifblank{#1}{\:\cdot\:}{#1}}
\DeclarePairedDelimiterXPP\normp[1]{}\lVert\rVert{_p}{\ifblank{#1}{\:\cdot\:}{#1}}
\DeclarePairedDelimiterXPP\normq[1]{}\lVert\rVert{_q}{\ifblank{#1}{\:\cdot\:}{#1}}
\DeclarePairedDelimiterXPP\norminf[1]{}\lVert\rVert{_\infty}{\ifblank{#1}{\:\cdot\:}{#1}}

%Bracket en angle
\DeclarePairedDelimiter\angl{\langle}{\rangle}

%Set, parenthèses
\newcommand{\set}[1]{\{ #1 \}}
\newcommand{\bigset}[1]{\bigl\{ #1 \bigr\}}
\newcommand{\Bigset}[1]{\Bigl\{ #1 \Bigr\}}
\newcommand{\bigpar}[1]{\bigl( #1 \bigr)}
\newcommand{\Bigpar}[1]{\Bigl( #1 \Bigr)}

%Opitmisation des opérateurs de comparaison
\DeclarePairedDelimiter\tio{o (}{)}
\DeclarePairedDelimiter\gro{O (}{)}

\renewcommand{\L} {\mathbb{L}}
\renewcommand{\F} {\mathbb{F}}
\newcommand{\U} {\mathcal{U}}
\newcommand{\D} {\mathcal{D}}

\pagestyle{fancy}

\fancyhead[C]{Consruction des corps finis}
\fancyhead[R]{}

\begin{document}

Dans ce développement, tout corps est supposé commutatif. Le symbole $\K$ désigne toujours un corps.

\section{Prérequis}

\subsection{Théorème de Lagrange}

Dans cette partie, on adopte une notation multiplicative pour les groupes. L'élément neutre est donc dénoté par $1$.

\begin{defn}[Ordre d'un groupe fini]
On appelle ordre d'un groupe fini son cardinal.
\end{defn}

\begin{thm}
Soit $(G, \cdot)$ un groupe et $H$ un sous groupe de $G$. La relation $R_g$ définie sur $G$ par :
\[ g_1 R g_2 \iff g_1^{-1}g_2 \in H \]
est une relation d'équivalence.
\end{thm}

\begin{proof}
Soit $g \in G$, on a $g^1 g = 1 \in H$, la relation est donc réflexive. \newline
Soient $g_1, g_2$ tels que $g_1 R g_2$, alors $(g_1^{-1}g_2)^{-1} = g_2^{-1}g_1 \in H$, la relation est donc symétrique. \newline
Soient $g_1, g_2, g_3$ tels que $g_1 R g_2$ et $g_2 R g_3$. Alors $g_1^{-1}g_2 g_2^{-1}g_3 = g_1^{-1}g_3 \in H$,
la relation est donc transitive.
\end{proof}

\begin{defn}[Classe à gauche]
Pour un groupe $G$ et $H$ un sous groupe de $G$ on appelle classe à gauche modulo $H$ d'un élément $g \in G$ l'ensemble 
\[ \bar{g} = \set{x \in G \tq g \R_g x} . \]
On remarquera que $x \in \bar{g} \iff g^{-1}x \in H \iff \exists h \in H \tq x = gh$. On a donc $\bar{g} = gH$.
\end{defn}

\begin{rmq}
On aurait pu définir la relation d'équivalence (et les classes) à droite.
\end{rmq}

%\begin{nota}
%On note $G/H$ l'ensemble des classes d'équivalence modulo $H$.
%\end{nota}

\begin{thm}
\label{thm:partGroup}
Soit $G$ un groupe et $H$ un sous-groupe de $G$, l'ensemble des classes à gauche modulo $H$ deux à deux distinctes
forment une partition de $G$.
\end{thm}

\begin{proof}
Notons $\bigcup_{i \in I} \widebar{g_i}$ l'ensemble des classes d'équivalence modulo $H$ deux à deux distinctes.
Soit $g \in G$, si $g \not \in \bigcup_{i \in I} \widebar{g_i}$ alors $\bar{g}$ est une classe d'équivalence distincte de toutes les classes
indicées par $I$ ce qui est exclu. Donc $G = \bigcup_{i \in I} \widebar{g_i}$. 
Supposons maintenant que $g \in \widebar{g_i} \cap \widebar{g_k}$ avec $i, k \in I$. 
Alors $g R g_i$ et $g R g_k$ donc, par symétrie et transitivité, $\widebar{g_i} = \widebar{g_k}$, et donc $i=k$.
\end{proof}

%\begin{defn}
%Pour un sous-groupe $H$ de $G$, le cardinal de $G/H$ est appelé l'indice de $H$ dans $G$ et on le note $[G : H]$.
%\end{defn}

\begin{thm}[Lagrange]
\label{thm:lagrange}
Soit $G$ un groupe fini d'ordre $n \geq 2$ et $H$ un sous-groupe de $G$. L'ordre de $H$ divise l'ordre de $G$.
\end{thm}

\begin{proof}
Soit $g \in G$ fixé, on considère l'application :
\[ \begin{array}{c c c c}
\varphi : & G & \to & G \\
& x & \mapsto & gx
\end{array} \]
qui est une permutation des éléments du groupe $G$. En conséquence, sa restriction à $H$ est une bijection de $H$ sur $gH$.
On conclut donc que $H$ et $gH$ ont même cardinal, et ce quelque soit $g \in G$. 
Comme l'ensemble des classes modulo $H$ deux à deux distinctes forment une partition de $G$ (cf. théorème $\ref{thm:partGroup}$, 
on en déduit que leur ordre (et donc celui de $H$) divise l'ordre de $G$. 
\end{proof}

%\subsection{Sous-groupe multiplicatif d'un corps} 
%
%\begin{defn}[Goupe monogène, cyclique]
%On dit qu'un groupe $G$ est monogène s'il existe un élément $g$ de $G$ tel que $G = \angl{g}$. 
%Si de plus $G$ est fini, on dit alors qu'il est cyclique.
%\end{defn}
%
%%\begin{rmq}
%%Pour un groupe multiplicatif cyclique, on a :
%%\[ \angl{g} = \Bigset{ \prod_{k=1}^r g^{\veps_k} ; \forall 1 \leq k \leq r, r \in \Ne, \veps_k \pm 1 } . \]
%%\end{rmq}
%
%\begin{defn}[Ordre d'un élément]
%L'ordre d'un élément $g$ d'un groupe $G$ est 
%\[ \theta(g) = \card{\angl{g}} \in \Ne \cup \set{\infty}. \]
%\end{defn}
%
%\begin{coro}
%Soit $G$ un groupe fini d'ordre $n \geq 1$. Alors l'ordre de tout élément $g$ de $G$ divise $n$.
%\end{coro}
%
%\begin{proof}
%il s'agit (pour $n \geq 2$) d'une conséquence du théorème $\ref{thm:lagrange}$.
%\end{proof}
%
%\begin{thm}
%\label{thm:ordppcm}
%Soit $(G, \cdot)$ un groupe abélien, $r \geq 2$ un entier et $g_1, \dots, g_r$ des éléments deux à deux distincts de $G$ 
%d'ordres respectifs $m_1, \dots, m_r \in \Ne$.
%Il existe dans $G$ un élément $g_0$ d'ordre égal au $\ppcm$ de ces ordres.
%\end{thm}
%
%\begin{proof}
%On utilisera le fait suivant :
%
%\begin{fait}
%Soit $a_1, a_2, \dots, a_n$ des nombres entiers, $n \geq 2$. On a :
%\[ \ppcm(\ppcm(a_1, \dots, a_{n-1}), a_n) = \ppcm(a_1, \dots, a_n) . \]
%\end{fait}
%
%On commence par le cas $r=2$. Si les ordres $m_1$ et $m_2$ sont premiers entre eux, alors $\ppcm(m_1, m_2) = m_1 m_2$ 
%et l'ordre de l'élément $g_0 = g_1 g_2$ est $m_1 m_2$. \newline
%Sinon, on peut toujours écrire :
%\[ m_1 = \prod_{i=1}^n p_i^{\alpha_i} \et m_2 = \prod_{i=1}^n p_i^{\beta_i}, \]
%avec les $p_i$ des nombres premiers et les $\alpha_i$ et $\beta_i$ des élements de $\N$ (éventuellement nuls).
%De plus, on organise ces décompositions en facteurs premiers de telle façon que :
%\[ \alpha_i > \beta_i \text{~pour~} i \leq k \text{~et~} \alpha_i \leq \beta_i \text{~sinon}, \]
%avec $0 \leq k \leq n+1$. On a alors :
%\[ m_1 = \prod_{i=1}^k p_i^{\alpha_i} \prod_{i=k+1}^n p_i^{\alpha_i} \et 
%	m_2 = \prod_{i=1}^k p_i^{\beta_i} \prod_{i=k+1}^n p_i^{\beta_i}, \]
%en fixant les produits vides à 1 si $k=0$ ou $k=n+1$. 
%On pose $q_1 =  \prod_{i=1}^k p_i^{\alpha_i}$ et $q_2 = \prod_{i=k+1}^n p_i^{\beta_i}$, 
%et on a $\ppcm(m_1, m_2) = q_1 q_2$, avec$q_1$ et $q_2$ premiers entre eux.
%Enfin, en écrivant $m_1 = q_1 r_1$ et $m_2 = q_2 r_2$, on trouve que les éléments $g_1^{r_1}$ et $g_2^{r_2}$ 
%sont d'ordres respectifs $q_1$ et $q_2$. 
%Donc, l'élément $g_0 = g_1^{r_1} g_2^{r_2}$ est d'ordre $q_1 q_2 = \ppcm(m_1, m_2)$. \newline
%
%Supposons le résultat vrai pour $r \geq 2$ quelconque et soient $g_1, \dots, g_{r+1}$ d'ordres respectifs $m_1, \dots, m_{r+1}$.
%Selon l'hypothèse de récurrence, il existe un élément $g_0'$ d'ordre $\ppcm(m_1, \dots, m_r)$ 
%et, selon le cas $r=2$, il existe un élément $g_0$ d'ordre :
%\[ \ppcm(\ppcm(m_1, \dots, m_r), m_{r+1}) = \ppcm(m_1, \dots, m_r, m_{r+1}). \]
%\end{proof}
%
%\begin{thm}
%\label{thm:ssgrpK}
%Soit $\K$ un corps. Tout sous-groupe fini du groupe multiplicatif $\Ket$ est cyclique.
%\end{thm}
%
%\begin{proof}
%On utilisera les faits suivants :
%
%\begin{fait}
%Dans un groupe fini, l'ordre de tous les éléments est fini.
%\end{fait}
%
%\begin{fait}
%Un polynôme de degré $n$ sur $\K[X]$ a, au plus, $n$ racines.
%\end{fait}
%
%\begin{fait}
%Un groupe fini d'ordre $n$ admettant un élément d'ordre $n$ est cyclique.
%\end{fait}
%
%Soit $(G, \cdot)$ un sous-groupe fini d'ordre $n$ de $\Ket$. Selon le théorème $\ref{thm:ordppcm}$,
%il existe dans $G$ un élément $g_0$ d'ordre $m \leq n$ égal au ppcm des ordres des éléments de $G$. 
%Ainsi, tout élément $g \in G$ est racine du polynôme $X^m-1$, donc $m \geq n$.
%On en conclut que $m=n$ et donc que $G$ est cyclique.
%\end{proof}

\subsection{Les anneaux $\Z$ et $\K[X]$}

\begin{defn}[Idéal]
Une partie $I$ d'un anneau $\A$ est appelée un idéal à gauche (resp. à droite) de $\A$ lorsque :
\begin{enumerate}[label=(\roman*)]
\item $(I,+)$ est un sous-groupe additif de $\A$ ;
\item pour tout $(x, a) \in I \times \A ; xa \in I$ (resp. $ax \in I$). 
\end{enumerate}
En particulier, un idéal est dit bilatère s'il est idéal à gauche et à droite.
\end{defn}

\begin{defn}[Idéal engendré]
Soit $P$ une partie d'un anneau $\A$, on appelle idéal engendré à gauche (resp. à droite) par $P$,
l'intersection de tous les idéaux à gauche (resp. à droite) de $\A$ contenant $P$.
\end{defn}

\begin{nota}
Pour $a \in \A$, on note $(a) = a \A = \{a \cdot q ; q\in \A\}$ l'idéal engendré par $a$.
\end{nota}

\begin{defn}[Idéal principal]
On dit qu'un idéal de $\A$ est principal s'il est engendré par un unique élément de $\A$.
\end{defn}

\begin{lem}[Sous-groupes de $\Z$]
\label{lem:nZ}
Les sous-groupes de $(\Z,+)$ sont les ensembles
\[ n\Z = \{ n \cdot k ; k \in \Z  \} \text{~pour~} n \in \N . \]
\end{lem}

\begin{proof}
Il est évident que $(n\Z, +)$ est un sous groupe de $(\Z,+)$. Par ailleurs, notons $G$ un sous-groupe de $\Z$.
Si $G$ ne contient qu'un seul élément, comme $0 \in G$ alors $G = 0\Z$. 
Sinon, $G$ contient au moins un élément non nul. Comme de plus $G$ est un groupe, il contient un élément strictement positif.
Notons $n = \min{ x \in G ; x > 0 }$, $n$ est bien défini et prenons $a$ un élément quelconque de $G$.
Si $a=0$, on écrit $a=0\cdot n$. Si $a >0$, en effectuant la division euclidienne de $a$ par $n$, on trouve :
\[ a = q \cdot n + r \text{~avec~} 0 \leq r < n . \]
Or, $r = a - q \cdot n \in G$ donc $r=0$ par définition de $n$, et $a = q \cdot n$.
Si $a<0$, on a $-a \in G$ et $-a = q \cdot n$ et donc $a=-q \cdot n$. 
On a donc $G \subset n\Z$. Il est de plus évident de $n\Z \subset G$.
Finalement, $\exists n \in \N \tq G = n\Z$. 
\end{proof}

\begin{thm}[Idéaux de $\Z$]
Les idéaux de $\Z$ sont les $n\Z$, pour $n \in \N$.
\end{thm}

\begin{proof}
Selon le lemme $\ref{lem:nZ}$, les idéaux de $\Z$, comme sous-groupes, sont nécessairement des éléments de $\{ n\Z ; n \in \N \}$.
Réciproquement, soit $n \in \N$ et soient $x \in n\Z, y \in \Z$, on a :
\[ x \cdot y = y \cdot x = n \cdot k \cdot y \in n\Z , \]
avec $k \in \Z \tq x = n \cdot k$.
\end{proof}

\begin{coro}
Tous les idéaux de $\Z$ sont principaux (et donc l'anneau $\Z$ est principal).
\end{coro}

\begin{thm}
\label{thm:ideauxKx}
Soit $\K$ un corps. L'anneau $\K[X]$ est principal. Précisément, pour tout idéal $I$ de l'anneau $\K[X]$, non réduit à $\set{0}$, 
il existe un unique polynôme unitaire $A$ tel que $I=(A)$. 
\end{thm}

\begin{proof}
Soit $I$ un idéal de $\K[X]$ non réduit à $\set{0}$. L'ensemble $\set{\deg{P} ; P \in I \setminus \set{0}}$ admet un minimum 
que l'on note $n_0$.
Prenons $A \in I \tq \deg{A} = n_0$ et soit $B \in I$, en effectuant la division euclidienne de $B$ par $A$ on obtient :
\[ B = AQ + R \text{~avec~} \deg{R} < n_0 ,\] 
donc, par définition de $A, R = 0$ et $I \subset (A)$.
Réciproquement, si $B \in (A), B = QA$ et donc $B \in I$ (car $I$ est un idéal).
Supposons maintenant que $(A)$ admette un second générateur $(B)$, alors $A$ divise $B$ et $B$ divise $A$, 
donc il existe $\lambda \in \Ket \tq A = \lambda B$. Le générateur unitaire est donc unique. \newline

Si $I= \set{0}$, alors $I=(0)$. Finalement, tous les idéaux de $\K[X]$ sont principaux.
\end{proof}

\subsection{Morphismes d'anneaux et idéaux}

\begin{defn}[Morphisme d'anneaux]
Soient $\A, \A'$ deux anneaux unitaires. 
Un morphisme d'anneaux et une application de $\A$ dans $\A'$, compatible avec les lois des anneaux,
et qui envoie le neutre multiplicatif de $\A$ sur le neutre multiplicatif de $\A'$.
\end{defn}

\begin{lem}
Soit $f$ un tel morphisme d'anneaux. Alors, pour tout $a \in \A$, on a
\begin{enumerate}[label=(\roman*)]
\item $f(0_{\A}) = 0_{\A'}$ ;
\item $f(-a) = -f(a)$.
\end{enumerate}
\end{lem}

\begin{proof}
\begin{enumerate}[label=(\roman*)]
\item Soit $a \in \A$, on a $f(a)=f(a+0_{\A})=f(a)+f(0_{\A})$ donc $f(0_{\A})=0_{\A'}$. 
\item Soit $a \in \A$, on a $f(0_{\A}) = f(a-a) = f(a) + f(-a) = 0_{\A'}$, d'où $f(-a) = -f(a)$.
\end{enumerate}
\end{proof}

\begin{thm}[Noyau et idéal]
\label{lem:noyau}
Soient $\A, \A'$ deux anneaux unitaires et $f$ un morphisme d'anneaux. 
Alors, le noyau de $f$ est un idéal (bilatère) de $\A$.
\end{thm}

\begin{proof}
Il est clair que $\ker(f) \subset \A$. Soient $x, y \in \ker(f)$, on a :
\[ f(x+y) = f(x) + f(y) = 0_{\A'} \et f(-x) = -f(x) = 0_{\A'} , \]
donc $\ker(f)$ est un sous-groupe additif de $\A$. \newline
Soit maintenant $a \in \ker(f)$, alors pour tout $b \in A, f(a \cdot b) = f(a) \cdot f(b) = 0_{\A'}$. De même pour $f(b \cdot a)$. 
Donc $\ker(f)$ est un idéal bilatère de $\A$.
\end{proof}

\begin{defn}
Soit $\A$ un anneau commutatif unitaire, soit $I$ un idéal de $\A$ et soient $a, b \in \A$. 
On dit que $a$ est congru à $b$ modulo $I$ si $b-a \in I$. 
On note alors $a\equiv b \pod{I}$ .
\end{defn}

\begin{thm}
La relation de congruence modulo $I$ est une relation d'équivalence sur $\A$.
\end{thm}

\begin{proof}
$I$ est un groupe additif donc $0 \in I$ donc $\forall a \in \A, a \equiv a \pod{I}$. 
Toujours de par la structure de groupe de $I, \forall a,b \in I$ si $b-a \in I$ alors $a-b \in I$.
La transitivité de la relation de congruence est évidente.
\end{proof}

\begin{nota}
Pour tout $a \in \A$, on note $\bar{a} = \{b \in \A ; b \equi a \pod{I} \}$ la classe d'équivalence de $a$ pour la relation de congruence.
On note $\A/I$ l'ensemble des classes d'équivalences modulo $I$.
\end{nota}

\begin{thm}
\label{thm:surjection}
Il existe une unique structure d'anneau commutatif unitaire sur $\A/I$ telle que la surjection canonique :
\[ \begin{array}{c c c c}
\pi_I : & \A & \to & \A/I \\
& a & \mapsto & \bar{a}
\end{array} \]
soit un morphisme d'anneaux.
\end{thm}

\begin{proof}
Soient $a,b \in I$ et $c \in \bar{a}, d \in \bar{b}$ on a :
\[ (c+d) - (a+b) = (c-a) + (d-b) \in I , \]
donc la relation de congruence est compatible avec l'addition. De même :
\[ cd - ab = c(d - b) - b(a - c) \in I , \]
donc la relation de congruence est aussi compatible avec la multiplication. 
On peut donc munir $\A/I$ d'une structure d'anneau avec les opérations ainsi définies
\[ \forall x, y \in \A/I \left \{ \begin{array}{c c c}
x + y & = & \widebar{a+b} \\
xy & = & \widebar{ab}
\end{array} \right. \]
avec $a$ et $b$ des éléments quelconques de $\bar{x}$ et $\bar{y}$ respectivement. 
En effet, si on prend  $a' \in \bar{x}$ et $b' \in \bar{y}$ distincts de $a$ et $b$, comme $a \equiv a'$ et $b \equiv b'$ on a :
\[ \widebar{a+b} = \widebar{a'+b'} \et \widebar{ab} = \widebar{a'b'} . \]
Les opérations ne dépendent donc pas des représentants choisis. 
On vérifie aussi facilement que les opérations définies sur $\A/I$ héritent des propriétés de commutativité (loi +), 
d'associativité et de commutativité (loi $\cdot$), et de distributivité (loi $\cdot$ par rapport à la loi $+$). \newline
On remarque que $\widebar{1_{\A}}$ est neutre pour la multiplication dans $\A/I$ donc $\A/I$ est unitaire.
Enfin, la compatibilité de la relation de congruence avec les lois $+$ et $\cdot$ démontre que
 $\pi_I$ est bien un morphisme d'anneaux, ce qui achève la démonstration de la structure mentionnée dans le théorème. \newline
 
Démontrons maintenant son unicité, si l'on suppose existante une telle structure, pour $a, b \in \A$, on a :
\[ \pi_I(a+b) = \pi_I(a) + \pi_I(b) \Rightarrow \widebar{a+b} = \bar{a} + \bar{b} , \]
et de même pour la multiplication. On revient donc à la structure précédemment exhibée. 
\end{proof}

%\begin{coro}
%Les idéaux de $\A$ sont les noyaux des morphismes d'anneaux $f : \A \to \A'$ où $\A'$ est un anneau commutatif unitaire.
%\end{coro}
%
%\begin{proof}
%On a déjà vu (cf. lemme $\ref{lem:noyau}$) que les noyaux de tels morphismes sont des idéaux de $\A$. 
%Réciproquement, si $I$ est un idéal de $\A$, alors $I$ est le noyau de la surjection de $\A$ dans $\A/I$ 
%définie dans le théorème $\ref{thm:surjection}$.
%\end{proof}

%\subsection{Idéal maximal}
%
%\begin{defn}
%Un idéal $I$ de $\A$ est dit maximal s'il est distinct de $\A$ et si $I$ et $A$ sont les seuls idéaux de $\A$ qui contiennent $I$.
%\end{defn}

%\begin{thm}
%\label{thm:idmax}
%Un idéal $I$ de $\A$ est maximal si, et seulement si, l'anneau quotient $\A/I$ est un corps.
%\end{thm}
%
%\begin{proof}
%
%Pour démontrer ce théorème, on utilisera le fait suivant :
%\begin{fait}
%Soient $I, J$ deux idéaux d'un anneau $\A$. Alors le sous-groupe $I+J = \{ x+y \tq x \in I, y \in J \}$ est un idéal de $\A$.
%\end{fait}
%
%Soit $I$ un idéal maximal de $\A$ et soit $a \not \in I$ (i.e. $\bar{a} \neq \bar{0}$).
%Le groupe $(a) + I$ est un idéal de $\A$ qui contient $I$, donc, par maximalité, $(a) + I = I \ou \A$.
%Dans le premier cas, $(a) = \{0\}$ et nécessairement $a=0$. Or ceci est exclu.
%Dans le second cas, il existe $(q, x) \in \A \times I \tq aq + x = 1$, soit $\bar{a}\bar{q} = \bar{1}$.
%Donc $\bar{a}$ est inversible dans $\A/I$, ce qui démontre que $\A/I$ est un corps. \newline
%
%Réciproquement, supposons que $\A/I$ est un corps et supposons que $I$ n'est pas maximal.
%Soit $J$ un idéal de $\A$ tel que $I \subsetneq J \subsetneq A$.
%Prenons $a \in J \setminus I$, on a $\bar{a} \neq 0$, il existe donc $(q,x) \in \A \times I$ tel que $aq + x = 1$.
%Or $aq \in J$ car $J$ est un idéal et $x \in I \subsetneq J$ donc $1 \in J$ et $J = \A$.
%On arrive à une contradiction, donc $I$ est maximal.
%\end{proof}

\subsection{Les anneaux $\Z/n\Z$}

\begin{thm}
L'ensemble $\Z/n\Z$, pour $n \in \N$ est un anneau commutatif unitaire. 
De plus, $\Z/n\Z = \{\widebar{0}, \widebar{1}, \dots, \widebar{n-1} \}$.
\end{thm}

\begin{proof}
Il s'agit d'une conséquence du théorème $\ref{thm:surjection}$.
\end{proof}

\begin{thm}
L'anneau $\Z/n\Z$ est un corps si, et seulement si, $n$ est un nombre premier.
\end{thm}

\begin{proof}
Soit $x \in \{1, \dots, n-1\}$, on cherche $u$ tel que :
\[ x \cdot u \equiv 1 \pod{n} \iff \exists v \in \Z \tq x \cdot u + n \cdot v = 1, \]
autrement dit, $x$ et $n$ sont premier entre eux.
Donc, $\Z/n\Z$ est intègre si $n$ est premier avec tous les entiers $\{1, \dots, n-1\}$. 
Ce qui revient à affirmer que $n$ est premier. \newline
Réciproquement, si $n$ est premier, il est premier avec tous les entiers $\{1, \dots, n-1\}$ ce qui implique que $\Z/n\Z$ est un corps.
\end{proof}

\begin{lem}
Tout anneau intègre fini est un corps.
\end{lem}

\begin{proof}
Soient $\A$ un anneau intègre fini et $a \in \A$, non nul. On considère l'application :
\[ \begin{array}{c c c c}
\sigma_a : & \A & \to & \A \\
& x & \mapsto & ax
\end{array} \]
Comme $\A$ est intègre, $\sigma_a$ est une injection. Comme l'anneau est fini, elle est bijective.
On peut donc définir sa réciproque et $\sigma_a^{-1}(1) = a^{-1}$. Ceci étant vrai pour tout $a \in \A$ non nul, $\A$ est bien un corps.
\end{proof}

\begin{coro}
\label{coro:integrZnZ}
L'anneau $\Z/n\Z$ est intègre si, et seulement si, $n$ est un nombre premier.
\end{coro}

\subsection{Les anneaux $\K[X]/(P)$}

\begin{thm}
L'ensemble $\K[X]/(P)$ pour $P \in \K[X]$ est un anneau commutatif unitaire. 
\end{thm}

\begin{proof}
Il s'agit d'une conséquence du théorème $\ref{thm:surjection}$.
\end{proof}

\begin{thm}
\label{thm:algbrKsP}
Soit $P \in \K[X]$ un polynôme unitaire de degré $n \geq 1$. 
\begin{enumerate}
\item L'algèbre $\K[X]/(P)$ est de dimension $n$ et $\bigpar{\widebar{X^k}}_{0 \leq k \leq n-1}$ en est une base.
\item Les conditions suivantes sont équivalentes :
	\begin{enumerate}[label=(\alph*)]
	\item $\K[X]/(P)$ est un corps ;
	\item l'anneau $\K[X]/(P)$ est intègre ;
	\item le polynôme $P$ est irréductible.
	\end{enumerate}
\end{enumerate}
\end{thm}


\begin{proof}
$\K[X]/(P)$ est un anneau selon le théorème $\ref{thm:surjection}$. De plus, la restriction de la surjection canonique de $\K[X]$ 
dans $\K[X]/(P)$ à $\K$ est injective ($a = 0 \cdot P +a$), on peut donc identifier tout scalaire $a$ à sa classe module $P$, 
ce qui permet de munir l'anneau quotient $\K[X]/(P)$ d'une structure de $\K$-espace vectoriel, et donc d'algèbre. \newline

\begin{enumerate}
\item
Soit $A \in \K[X]$, par division euclidienne, $A = PQ + R$ avec $R \in \K_{n-1}[X]$ et 
$\bar{A} = \bar{P} = \sum_{n=0}^{k-1} \widebar{X^k}$, donc $\bigpar{\widebar{X^k}}_{0 \leq k \leq n-1}$ est une famille génératrice 
de $\K[X]/(P)$.
De plus :
\[ \bar{R} = \sum_{k=0}^{n-1} a_ k \widebar{X^k} = \bar{0} \iff R \in \bar{P} . \] 
Comme le degré de $R$ est strictement inférieur à celui de $P$, ceci est équivalent à :
\[ R = 0 \iff \forall k \in \set{0, \dots, n-1}, a_k=0 . \]
La famille $\bigpar{\widebar{X^k}}_{0 \leq k \leq n-1}$ est donc une base de $\K[X]/(P)$. \newline

\item 
Si $\K[X]/(P)$ est un corps, alors c'est un anneau intègre. Si c'est un anneau intègre, une égalité $P=AB$, 
avec $A, B$ de degré au moins 1 donnerait $\bar{A}\bar{B} = \bar{0}$ ce qui est contradictoire.
Si le polynôme $P$ est irréductible, alors tout polynôme non nul $A$ de degré strictement inférieur à $n$ est premier avec $P$ et,
selon le théorème de Bézout, il existe deux polynômes $U$ et $V$ tels que $AU+PV=1$, soit $\bar{A}\bar{U} = 1$.
Or, tout élément de $\K[X]/(P)$ est de degré au plus $n-1$ selon le point 1., donc inversible.
\end{enumerate}
\end{proof}

\begin{rmq}
On remarquera que, par compatibilité de la loi $\cdot$ dans $\K[X]/(P)$, on a $\widebar{X^k} = \widebar{X}^k$.
\end{rmq}

\subsection{Caractéristique d'un corps}

\begin{thm}
\label{thm:morpCarac}
On suppose que $\A$ est un anneau unitaire. Alors l'application
\[ \begin{array}{c c c c}
\varphi : & \Z & \to & \A \\
& n & \mapsto & n \cdot 1_{\A}
\end{array} \]
est l'unique morphisme d'anneau de $\Z$ dans $\A$.
\end{thm}

\begin{proof}
$\varphi$ est bien un morphisme d'anneaux par associativité et distribituvité de la loi $\cdot$ par rapport à la loi $+$. \newline
Supposons qu'il existe un second tel morphisme que l'on note $f$. 
Soit $n \in \N$, on peut écrire $n = \underbrace{1 + \dots + 1}_{\text{n fois}}$ et donc $f(n) = n.1_{\A}$. 
Pour $n \tq -n \in \N$, on raisonne de la même manière en remarquant que $f(-1) = -1_{\A}$.
\end{proof}

\begin{rmq}
Comme $\varphi$ est un morphisme d'anneau, d'après le lemme $\ref{lem:noyau}$, son noyau est un idéal de $\Z$,
il existe donc $p \in \N$ tel que $\ker(\varphi)=p\Z$.
\end{rmq}

\begin{defn}[Caractéristique d'un anneau unitaire]
Soit $\A$ un anneau unitaire, on définit la caractéristique de $\A$ comme l'entier naturel $p$ 
tel que $p\Z$ est le noyau du morphisme $\varphi$.
\end{defn}

\begin{lem}
\label{lem:caracSS}
Soit $\A'$ un sous-anneau d'un anneau unitaire $\A$, alors la caractéristique de $\A'$ est égale à la caractéristique de $\A$.
\end{lem}

\begin{proof}
Il suffit de restreindre $\varphi$ à $\A'$ et de considérer le fait que $1_{\A} = 1_{\A'}$.
\end{proof}

\begin{lem}
\label{lem:isoKer}
L'application :
\[ \begin{array}{c c c c c}
\varphi' : & \Z/\ker{\varphi} & \to & \A & \\
& x & \mapsto & \varphi(a) & \text{~avec~} a \in x
\end{array} \]
est un morphisme d'anneaux injectif.
\end{lem}

\begin{proof}
On démontre d'abord que $\varphi'$ est bien définie, c'est-à-dire que l'image de $x \in \Z/\ker{\varphi}$ par $\varphi'$ 
ne dépend pas du représentant $a \in x$ choisi.
Soient $a, a' \in x$, il existe $k \in \Z \tq a'=a+k \cdot p$ et donc :
\[ \varphi(a') = \varphi(a) + \varphi(k) \cdot \varphi(p) = \varphi(a) , \]
car $\ker(\varphi) = p\Z$.
Ainsi $\varphi'$ est bien défini, c'est donc un morphisme d'anneaux car $\varphi$ est un morphisme d'anneaux.
Il est injectif par construction.
\end{proof}

\begin{thm}[Cas d'un anneau unitaire intègre]
\label{thm:caracIntegre}
La caractéristique d'un anneau unitaire intègre $\A$ est 0 ou un nombre premier. 
Dans le cas où la caractéristique est nulle, l'anneau est infini.
\end{thm}

\begin{proof}
Pour démontrer ce théorème, on utilise les deux faits suivants :

\begin{fait}
Soit $f: \A \to \A'$ un morphisme d'anneaux. Alors $f(\A)$ est un sous-anneau de $\A'$.
\end{fait}

\begin{fait}
Soit $\A$ un anneau intègre et $\A'$ un sous-anneau de $\A$. Alors $\A'$ est intègre.
\end{fait}

Supposons que $p \neq 0$, selon le lemme $\ref{lem:isoKer}$ l'anneau $\Z/p\Z$ est ismorphe à son image par $\varphi'$,
c'est-à-dire à un sous-anneau de $\A$. Donc $\Z/p\Z$ est isomorphe à un anneau intègre, il est donc intègre 
ce qui implique que $p$ est premier (cf corollaire $\ref{coro:integrZnZ}$) \newline
Si $p=0$ alors $\varphi$ est un morphisme d'anneaux injectif, 
donc $\varphi(\Z)$ est un sous-anneau de $\A$, isomorphe à $\Z$ donc infini, ce qui implique que $\A$ est infini.
\end{proof}

\begin{coro}
La caractéristique d'un corps est 0 ou un nombre premier.
\end{coro}

\subsection{Sous-corps premier}

Dans cette partie, pour un nombre $p$ premier, on écrira $\F_p = \Z / p\Z$.

\begin{defn}
Le plus petit sous-corps d'un corps $\K$, noté $\K_0$ est appelé sous-corps premier de $\K$.
\end{defn}

\begin{lem}
\label{lem:kev}
Soit $\L$ un corps et $\K$ un sous-corps de $\L$. Alors $\L$ est muni d'une structure de $\K$-espace vectoriel.
Si de plus $\L$ est fini, alors il existe $n \in \N$ tel que $\card{~\L} = \card{~\K}^n$
\end{lem}

\begin{proof}
$(\L, +, \cdot)$ est un corps, donc $(\L, +)$ est un groupe abélien.  De plus, en restreignant la loi $\cdot$ sur $(\K \times \L)$, 
on aura bien une loi de composition externe distributive et vérifiant une associativité mixte avec la loi $+$.
Enfin, l'ensemble des éléments de $\L$ est une famille génératrice du $\K$-e.v. $\L$, ce qui implique que sa dimension est finie.
\end{proof}

\begin{lem}
\label{lem:isoFp}
Si $\K$ est de caractéristique $p \geq 2$ premier, le sous-corps premier de $\K$ est isomorphe à $\F_p$.
\end{lem}

\begin{proof}
On utilisera le fait suivant :

\begin{fait}
Tout sous-corps fini d'un corps de caractéristique $p \geq 2$ est de cardinal au moins $p$.
\end{fait}

On considère le morphisme $\varphi : \Z \to \K$ introduit au théorème $\ref{thm:morpCarac}$. 
Par définition, son noyau est $p \Z$ et donc le corps $\Z / p \Z$ est isomorphe à son image par $\varphi$ que l'on note $\K_0$.
Or $(\K_0, +)$ est un groupe d'ordre $p$ premier donc, selon le théorème de Lagrange, les sous-groupes de $(\K_0, +)$ sont $\K_0$
et $\set{0}$. A fortiori, le seul sous-corps de $\K_0$ est $\K_0$. C'est donc le sous-corps premier de $\K$.
\end{proof}

%\begin{lem}
%$\A$ est un corps si, et seulement si, les seuls idéaux de $\A$ sont $\A$ et $\{0\}$.
%\end{lem}
%
%\begin{proof}
%Supposons que $\A$ est un corps. Considérons $I$ un idéal de $\A$ que l'on suppose distinct de $\A$.
%Prenons $a \in I \setminus \{0\}$, comme $a$ est inversible, il existe $b \in \A \tq ab = 1$.
%Donc $1 \in I$ et $I = \A$ ce qui est exclu. Donc $I = \{0\}$. \newline
%
%Supposons maintenant que les seuls idéaux de $\A$ soient $\A$ et $\{0\}$.
%Soit $a \in \A$ non nul, alors $(a) = \A$, donc il existe $b \in \A \tq ab = 1$. $\A$ est bien un corps.
%\end{proof}

%\begin{lem}
%\label{lem:injectif}
%Un morphisme d'anneaux de $\K$ dans un anneau unitaire $\A$ est nécessairement injectif.
%\end{lem}
%
%\begin{proof}
%Le noyau d'un tel morphisme est un idéal du corps $\K$, c'est donc tout le corps ou $\{0\}$. 
%Mais comme le morphisme envoie $1_{\K}$ sur $1_{\A}$, son noyau est nécessairement réduit à $\{0\}$. 
%\end{proof}

%\begin{defn}[Extension de corps]
%Une extension d'un corps $\K$ est la donnée d'un couple $(\L, \sigma)$ formé d'un corps commutatif $\L$ 
%et d'un morphisme de corps $\sigma : \K \to \L$.
%\end{defn}
%
%\begin{nota}
%Pour tout $\omega$ dans une extension $\L$ de $\K$, on note $\K[\omega] = \bigset{P(\omega) ; P \in \K[X]}$,
%et on désigne par $\K(\omega)$ le plus petit sous-corps de$\L$ qui contient $\K$ et $\omega$.
%\end{nota}
%
%\begin{defn}
%On dit qu'un élément $\omega$ d'une extension $\L$ de $\K$ est algébrique sur $\K$, s'il existe un polynôme non nul $P$ 
%dans $\K[X]$ tel que $P(\omega)=0$.
%\end{defn}
%
%\begin{lem}[Polynôme minmal]
%Si un élément $\omega$ d'une extension $\L$ de $\K$ est algébrique sur $\K$, il existe un unique polynôme unitaire $P_{\omega}$
%tel que $P_{\omega}(\omega)=0$. Plus précisément, $\bigset{P \in \K[X] ; P(\omega) = 0} = (P_{\omega})$.
%\end{lem}
%
%\begin{proof}
%On considère le morphisme d'anneaux de $\K[X]$ sur $\L$, son noyau est un idéal de $\K[X]$, d'où le résultat 
%(cf. théorème $\ref{thm:ideauxKx}$).
%\end{proof}
%
%\begin{defn}[Corps de rupture]
%On dit qu'une extension $\L$ de $\K$ est un corps de rupture d'un polynôme non constant $P \in \K[X]$, si le polynôme $P$
%a une racine $\omega$ dans $\L$ telle que $\L=\K[\omega]$.
%\end{defn}
%
%\begin{thm}
%Si $P(X) = \sum_{k=0}^n a_k X^k$ est un polynôme irréductible unitaire de degré $n \geq 1$ dans $\K[X]$, 
%alors $\K[X]/(P)$ est un corps de rupture de $P$ et $P$ est le polynôme minimal de $\omega = \widebar{X}$ sur $\K$.
%\end{thm}
%
%\begin{proof}
%Selon le théorème $\ref{thm:algbrKsP}$, le couple $\bigpar{\K[X]/(P), \sigma}$ (en posant $\sigma : x \in \K \mapsto \bar{x}$) 
%est une extension de corps sur $\K$ de degré $n$ et de base 
%\[ \bigpar{\widebar{X^k}}_{0 \leq k \leq n-1} = (\widebar{\omega^k})_{0 \leq k \leq n-1} . \]
%De plus, pour tout $Q = \sum_{k=0}^n b_k X^k \in \K[X]$
%\[ Q(\omega) = \sum_{k=0}^n b_k \omega^k = \sum_{k=0}^n b_k \widebar{X}^k = \widebar{Q}, \]
%en particulier $P(\omega) = \widebar{P} = 0$. 
%Finalement, $\K[X]/(P) = \K[\omega]$, ce qui démontre que $\K[X]/(P)$ est un corps de rupture de $P$. \\
%Enfin, pour $Q \in \K[X] $ on a $(Q(\omega)=\bar{0}) \iff (\bar{Q} = \bar{0} \iff Q \in (P)$. 
%$P$ étant unitaire, il est le polynôme minimal de $\omega$. 
%\end{proof}

\section{Construction de corps finis}

\begin{limi}
Dans tout ce paragraphe, $p \geq 2$ est un nombre premier. \newline

Pour tout entier $n \geq 1$, on note $\U_n(p)$ l'ensemble des polynômes unitaires irréductibles de degré $n$ dans $\F_p[X]$
et $I_n(p)$ le cardinal de $\U_n(p)$. On note aussi $\D_n$ l'ensemble des diviseurs de $n$ dans $\Ne$. \newline

Pour tout polynôme $P \in \U_n(p)$, le quotient $\F_p[X]/(P)$ est une $\F_p$-algèbre de dimension $n$, 
de base $\bigpar{\widebar{X^k}}_{0 \leq k \leq n-1}$ et c'est donc un corps fini de cardinal $p^n$ 
(cf. théorème $\ref{thm:algbrKsP}$). \newline

Pour $P \in \U_n(p)$, on notera $\F_{p^n} = \F_p[X]/(P)$.
\end{limi}

\begin{thm}
\label{thm:cardinal}
Si $\K$ est un corps fini, il est de cardinal $p^n$, où $p \geq 2$ est un nombre premier (sa caractéristique) et $n \in \Ne$.
Dans ce cas, tout sous-corps de $\K$ est de cardinal $p^d$ où $d$ est un diviseur de $n$.
Réciproquement, pour tout diviseur $d$ de $n$, il existe un unique sous-corps de $\K$ de cardinal $p^d$,
à savoir le corps $\F = \Bigset{x \in \F_{p^n} ; x^{p^d} = x }$.
\end{thm}

\begin{proof}

Selon le théorème $\ref{thm:caracIntegre}$, la caractéristique du corps $\K$ est un nombre premier que l'on note $p$.
Par ailleurs, $\K$ peut être muni d'une structure de $\K_0$-espace vectoriel de dimension finie (cf. lemme $\ref{lem:kev}$),
donc $\K$ est isomorphe à $(\K_0)^n$ avec $n \in \Ne$. Or, selon le lemme $\ref{lem:isoFp}, \K_0$ est isomorphe à $\F_p$.
Finalement, $\card{K} = \bigpar{\card{F_p}}^n = p^n$. \newline

Considérons maintenant un sous-corps $\F$ de $\K$. 
La caractéristique de $\F$ est la même que celle de $\K$ (cf. lemme $\ref{lem:caracSS}$, on a donc $\card{F} = p^d$ avec $d \leq n$. 
Or, le corps $\K$ pouvant être muni d'un structure de $\F$-espace vectoriel, il est isomorphe à
$\F^q$ avec $q \in \Ne$, donc $p^n = (p^d)^q$ soit $n = dq$, donc $d$ est bien un diviseur de $n$. \newline

Réciproquement, soit $d$ un diviseur de $n$ et soit $\F$ un éventuel sous-corps de $\K$ de cardinal $p^d$. 
Soit $x \in \F$, selon le théorème de Lagrange, le groupe monogène $\angl{x} = \bigset{x^n ; n \in \N}$ 
est d'ordre divisant l'ordre de $\Ket$ = $p^d-1$.
Il suit que, pour tout $x \in \F, x^{p^d} = x$ et donc que le sous-corps $\F$ est inclus dans l'ensemble des racines du polynôme 
$P_d(X) = X^{p^d} - X$. \\
Comme ce polynôme a, au plus, $p^d$ racine et que $\card{F} = p^d$, on conclut que, si un tel $\F$ existe, il est unique et égal à
l'ensemble des racines du polynôme $P_d(X)$. \newline

Il reste à montrer que l'ensemble des racines du polynôme $P_d(X)$ (noté $\F$) est bien un sous-corps de $\K$ à $p^d$ éléments.
Il est clair que $P_d(0) =P_d(1) = 0$. Soient $x, y \in \F$, on a :
\[ (x-y)^{p^d} = x^{p^d} + \sum_{k=1}^{p^d-1} \binom{p^d}{k} x^k (-y)^{p^d-k} + (-y)^{p^d}, \]
or, $p$ étant premier il divise tous les coefficients binomiaux 
\[ \binom{p^d}{k}, k \in \set{1, \dots, p^d-1} , \]
de plus, $p$ étant soit égal à 2 (et dans ce cas $1=-1$) soit impair, on a :
\[ (x-y)^{p^d} = x^{p^d} - y^{p^d} = x - y, \]
Donc $(\F, +)$ est un groupe abélien. 
Soient $x, y \in \F^{\ast}$, on a 
\[ \Bigpar{\dfrac{x}{y}}^{p^d} = \dfrac{x^{p^d}}{y^{p^d}} = \dfrac{x}{y}, \]
Donc $(\F^{\ast}, \cdot)$ est un groupe abélien. \\
Ensuite, comme $d$ divise $n$, il existe $q \in \Z$ tel que $n = qd$ et 
\[ p^n-1 = p^{qd}-1 = (p^d-1)(1+p^d+\dots+p^{(q-1)d}) = (p^d-1)m . \]
Et donc :
\[ X^{p^n-1} - 1 = X^{(p^d-1)m}-1 = (X^{p^d-1} -1)(1 + X^{p^d-1} + \dots + X^{(p^d-1)(m-1)}) . \]
Comme $X^{p^n-1}-1$ admet tous les éléments de $\Ket$ pour racines (toujours selon le théorème de Lagrange), 
il en est donc de même pour $X^{p^d-1} -1$ et finalement $P_d(X) = (X^{p^d-1} -1)X$ est scindé avec $p^d$ racines dans $\K$.
\end{proof}

\begin{lem}
\label{lem:PdivPn}
Pour tout diviseur $d$ de $n$, tout polynôme $P \in \U_d(p)$ divise le polynôme $P_n(X) = X^{p^n} - X$ dans $\F_p[X]$.
\end{lem}

\begin{proof}
Soit $P \in \U_d(p), \F_{p^d}$ est un corps de cardinal $p^d$, donc $\F_{p^d}^{\ast}$ est un corps multiplicatif de cardinal $p^d-1$,
et $\bar{X}^{p^d-1} = \bar{1}$ (théorème de Lagrange), donc $\bar{X}^{p^d} = \bar{X}$. Par récurrence, on peut montrer que 
$\bar{X}^{p^{kd}} = \bar{X}$ pour tout $k \in \N$ (en effet $\bar{X}^{p^{(k+1)d}} = \bigpar { \bar{X}^{p^{kd}} }^{p^d}$). \\
Comme $d$ divise $n$, il existe $q$ tel que $n=dq$ et par conséquence, $\bar{X}^{p^n} = \bar{X}^{p^{qd}} = \bar{X}$, 
autrement dit, $\widebar{P_n} = \bar{0}$ dans $F_{p^d}$.
\end{proof}

\begin{thm}[Unicité des corps finis]
A un isomorphisme près, il n'existe qu'un seul corps à $p^n$ éléments, c'est le corps $\F_{p^n} = \F_p[X] / (P)$,
où $P \in U_n(p)$.
\end{thm}

\begin{proof}
On désigne par $\K$ un corps fini à $p^n$ éléments, donc de caractéristique $p$ (cf. théorème $\ref{thm:cardinal}$).
Selon le lemme $\ref{lem:isoFp}$, le sous-corps premier de $\K$ est isomoprhe à $\F_p$. 
On peut donc identifier tout polynôme dans $\F_p[X]$ à un polynôme dans $\K[X]$. \newline

Selon le lemme $\ref{lem:PdivPn}$, le polynôme $P$ divise $P_n(X) = X^{p^n} - X$ dans $F_p[X]$ donc dans $\K[X]$. 
Or, toujours selon le théorème de Lagrange, comme $\card{\Ket} = p^n-1$, on a $\lambda^{p^n} = \lambda$ 
pour tout $\lambda$ dans $\K$ donc  le polynôme $P_n$ est nul pour tout élément de $\K$, on peut donc écrire :
\[ P_n(X) = \prod_{\lambda \in \K} (X - \lambda) . \]
$P_n$ est scindé et à racines simples dans $\K[X]$, il en est de même pour $P$ qui divise $P_n$.
Considérons maintenant l'application :
\[ \begin{array}{c c c c}
f : & \F_p[X] & \to & \K \\
& Q & \mapsto & Q(\lambda) , 
\end{array} \]
avec $\lambda \in \K$ racine de $P$. $f$ est un morphisme d'anneaux et son noyau est donc un idéal de $\K[X]$.
Comme l'anneau $\K[X]$ est principal (cf. théorème $\ref{thm:ideauxKx}$), il existe $Q \in \F_p[X]$ unitaire et tel que $\ker{f} = (Q)$.
De plus, l'idéal $(P)$ est inclus dans $(Q)$, donc le polynôme $Q$ divise le polynôme unitaire irréductible $P$ dans $\F_p[X]$.
Le cas $Q=1$ est exclu, car alors $\lambda$ est racine de tout polynôme de $\F_p[X]$ ce qui est absurde. 
On conclut donc que $P=Q$ et, par suite, $\ker{f} = (P)$. La restriction de $f$ à $\F_{p^n}$ est donc une injection dans $\K$.
Ces deux corps étant de même cardinalité (voir limaire), on conclut qu'ils sont isomorphes.
\end{proof}

\end{document}